\section{Experimental Protocol}
\label{sec:protocol}

\subsection{Frozen Execution and Provenance}

Scientific workers ran in Kaggle sessions with one visible NVIDIA T4; the device
mask was set before importing CUDA-aware packages. The initial stage froze source
and configuration hashes, route specifications, prompts, model and encoder
revisions, decoding, seed 4622, and the Python/Torch/CUDA environment. Later stages
rejected a changed experiment identifier, specification hash, source bundle,
environment contract, or worker lane. Deterministic decoding used FP16, greedy
generation, an 80-token output limit, and a 1,536-token context budget. OOM or
context overflow was recorded rather than silently changing precision or budget.

The pilot evaluated 22 actions on 80 queries and pruned five non-anchor actions by
the frozen rule. Router fitting used 11 eligible actions over 400 train, 120
calibration, and 80 validation queries, yielding 4,400, 1,320, and 880 route--query
traces. The test matrix was not used for fitting, calibration, thresholding, or
fallback selection. The clean release executes each of 17 retained actions on all
120 test queries once, for 2,040 traces.

\subsection{Two-Stage Controller}

Stage 0 uses eight query-available features---length, character and digit counts,
entity codes, and temporal, conflict, memory, and multi-hop cues---to score direct
actions. If no direct action clears the decision constraints, one charged
BM25/memory-metadata probe supplies eight retrieval statistics, including top
scores, score gaps, counts, and authority values. Stage 1 then scores non-direct
actions. If no action qualifies, the controller invokes a validation-selected
fail-closed fallback.

For each of five grouped-bootstrap seeds, a histogram gradient-boosting classifier
predicts strict success. Isotonic regression calibrates each predictor on the
calibration split, following the general requirement that decision confidence be
empirically calibrated \cite{GuoEtAl2017Calibration}. Separate quantile regressors
predict log energy and latency at $q=0.90$. The controller uses the minimum
calibrated success probability across bootstrap members and mean exponentiated cost
predictions. Among routes above threshold $\tau$ with predicted latency at most
12 s, it chooses the lowest predicted generation energy. This reject-option design
is related to selective prediction \cite{GeifmanElYaniv2019SelectiveNet}, but its
fallback is not evidence of safety; it is only the declared completion behavior.

\subsection{Frozen Constraints and Amendment}
\label{sec:amendment}

A validation threshold was feasible only if strict success was at least 0.80,
execution coverage at least 0.90, and abstention at most 0.20. The 80\% value was a
predeclared deployment utility requirement, not a prediction of attainable model
performance. It was intentionally not relaxed after observing the routes. Every
threshold in the original grid $\{0.60,0.65,\ldots,0.85\}$ produced 8.75\%
validation success, 100\% coverage, and 18.75\% abstention, so the validation
constraints were infeasible and the original implementation stopped without a
policy artifact.

After validation infeasibility was known but before test-label access, amendment
\modelid{stage06-threshold-completion-separation-v1} extended the audit grid to
0.90, 0.95, and 1.00. Outcomes were unchanged. The amendment froze $\tau=1.0$ to
permit completion of the already specified test matrix, recorded
\texttt{validation\_feasible=false}, and prohibited a positive hypothesis claim.
The fallback A03 was selected on validation by strict success, then median measured
GPU energy, then route identifier. We retain the amendment as a protocol deviation
and never describe this fallback as validation-feasible or empirically safe.

\subsection{Confirmatory Estimands}

For test query $i$, let $\pi(i)$ be the frozen policy action and $b$ the fixed A03
fallback baseline. The paired estimands are

\begin{align}
\Delta_S &= \frac{1}{N}\sum_{i=1}^{N}
\left(S_{i,\pi(i)}-S_{i,b}\right), \\
\Delta_E &= \frac{1}{N}\sum_{i=1}^{N}
\left(E_{i,\pi(i)}-E_{i,b}\right).
\end{align}

The confirmatory claim required all four gates:

\begin{equation}
\begin{aligned}
L_{.95}(\Delta_S)&\geq-0.03, & U_{.95}(\Delta_E)&<0, \\
\mathrm{coverage}&\geq0.90, & \mathrm{abstention}&\leq0.20.
\end{aligned}
\label{eq:gates}
\end{equation}

If policy and baseline select the same trace, both paired effects are identities at
zero. Passing the formal non-inferiority margin in that case does not demonstrate
preservation of useful quality; energy reduction still requires a strictly
negative upper interval bound.

\subsection{Uncertainty and Diagnostic Analyses}

Paired 95\% intervals use 10,000 percentile-bootstrap replicates with seed 4622,
resampling whole scenario IDs \cite{EfronTibshirani1994Bootstrap}. They quantify
variation across the frozen questions, not decoding or repeated-hardware variance.
The predeclared analyses cover policy--baseline effects, matched
grounded-minus-direct effects, and route-level contrasts.

After the confirmatory failure, we derived the action-pool oracle, routing
headroom, unique successes, generator interaction contrasts, task-aware oracle,
retrieval-to-utilization decomposition, and answer-F1 threshold sensitivity from
the same immutable traces. These are explicitly exploratory and are not
multiplicity-adjusted. No model was rerun and no outcome was replaced.

\subsection{Robustness and Hardware Assignment}

The diagnostic corruption matrix evaluates A00 direct, A03 hybrid, and A13
grounded/guarded under stale evidence, untrusted conflict injection, missing
evidence, and deletion/duplicate ordering, for 1,440 traces. Because the selected
policy has zero clean utility, its corruption deltas are interpreted only as a
floor effect. The literal compromise marker is a narrow synthetic indicator, not a
general injection detector.

Clean routes span four physical T4 boards. Four matched pairs---Qwen3-0.6B,
Granite-4.1-3B, SmolLM3-3B, and Qwen3-4B---are same-board within pair. The
Granite-4.0-1B pair is cross-board. Tiny BM25, dense, hybrid, and direct endpoints
share one board; the tiny memory routes do not. The policy and A03 baseline refer
to the same stored trace on every query, so their identity comparison is unaffected
by board assignment.
